\begin{filecontents*}{neurips_2026.sty}
\NeedsTeXFormat{LaTeX2e}
\ProvidesPackage{neurips_2026}[2026/05/01 v1.0 NeurIPS 2026 submission/camera-ready style file]

\newif\if@neuripsfinal\@neuripsfinalfalse
\DeclareOption{final}{\@neuripsfinaltrue}
\newif\if@natbib\@natbibtrue
\DeclareOption{nonatbib}{\@natbibfalse}
\newif\if@preprint\@preprintfalse
\DeclareOption{preprint}{\@preprinttrue}
\DeclareOption{main}{}
\DeclareOption{position}{}
\DeclareOption{eandd}{}
\DeclareOption{creativeai}{}
\DeclareOption{sglblindworkshop}{}
\DeclareOption{dblblindworkshop}{}
\DeclareOption{education}{}
\ProcessOptions\relax

\newif\if@submission\@submissiontrue
\if@neuripsfinal\@submissionfalse\fi
\if@preprint\@submissionfalse\fi

\renewcommand{\rmdefault}{ptm}
\renewcommand{\sfdefault}{phv}

\newcommand{\@neuripsordinal}{40th}
\newcommand{\@neuripsyear}{2026}
\newcommand{\@neuripslocation}{Sydney, Australia}

\newcommand{\acksection}{\section*{Acknowledgments and Disclosure of Funding}}
\newbox\@neuripsackbox
\newenvironment{ack}{%
  \if@submission
    \setbox\@neuripsackbox\vbox\bgroup
  \else
    \acksection
  \fi
}{%
  \if@submission
    \egroup
  \fi
}

\if@natbib
  \RequirePackage{natbib}
\fi

\RequirePackage[verbose=true,letterpaper]{geometry}
\AtBeginDocument{
  \newgeometry{
    textheight=9in,
    textwidth=5.5in,
    top=1in,
    headheight=12pt,
    headsep=25pt,
    footskip=30pt
  }
}

\widowpenalty=10000
\clubpenalty=10000
\flushbottom
\sloppy

\renewcommand{\normalsize}{%
  \@setfontsize\normalsize\@xpt\@xipt
  \abovedisplayskip      7\p@ \@plus 2\p@ \@minus 5\p@
  \abovedisplayshortskip \z@ \@plus 3\p@
  \belowdisplayskip      \abovedisplayskip
  \belowdisplayshortskip 4\p@ \@plus 3\p@ \@minus 3\p@
}
\normalsize
\renewcommand{\small}{%
  \@setfontsize\small\@ixpt\@xpt
  \abovedisplayskip      6\p@ \@plus 1.5\p@ \@minus 4\p@
  \abovedisplayshortskip \z@  \@plus 2\p@
  \belowdisplayskip      \abovedisplayskip
  \belowdisplayshortskip 3\p@ \@plus 2\p@   \@minus 2\p@
}
\renewcommand{\footnotesize}{\@setfontsize\footnotesize\@ixpt\@xpt}
\renewcommand{\scriptsize}{\@setfontsize\scriptsize\@viipt\@viiipt}
\renewcommand{\tiny}{\@setfontsize\tiny\@vipt\@viipt}
\renewcommand{\large}{\@setfontsize\large\@xiipt{14}}
\renewcommand{\Large}{\@setfontsize\Large\@xivpt{16}}
\renewcommand{\LARGE}{\@setfontsize\LARGE\@xviipt{20}}
\renewcommand{\huge}{\@setfontsize\huge\@xxpt{23}}
\renewcommand{\Huge}{\@setfontsize\Huge\@xxvpt{28}}

\providecommand{\section}{}
\renewcommand{\section}{%
  \@startsection{section}{1}{\z@}%
                {-2.0ex \@plus -0.5ex \@minus -0.2ex}%
                { 1.5ex \@plus  0.3ex \@minus  0.2ex}%
                {\large\bf\raggedright}%
}
\providecommand{\subsection}{}
\renewcommand{\subsection}{%
  \@startsection{subsection}{2}{\z@}%
                {-1.8ex \@plus -0.5ex \@minus -0.2ex}%
                { 0.8ex \@plus  0.2ex}%
                {\normalsize\bf\raggedright}%
}
\providecommand{\subsubsection}{}
\renewcommand{\subsubsection}{%
  \@startsection{subsubsection}{3}{\z@}%
                {-1.5ex \@plus -0.5ex \@minus -0.2ex}%
                { 0.5ex \@plus  0.2ex}%
                {\normalsize\bf\raggedright}%
}
\providecommand{\paragraph}{}
\renewcommand{\paragraph}{%
  \@startsection{paragraph}{4}{\z@}%
                {1.5ex \@plus 0.5ex \@minus 0.2ex}%
                {-1em}%
                {\normalsize\bf}%
}

\setlength{\parindent}{\z@}
\setlength{\parskip}{5.5\p@}

\setlength{\topsep}{4\p@ \@plus 1\p@ \@minus 2\p@}
\setlength{\partopsep}{1\p@ \@plus 0.5\p@ \@minus 0.5\p@}
\setlength{\itemsep}{2\p@ \@plus 1\p@ \@minus 0.5\p@}
\setlength{\parsep}{2\p@ \@plus 1\p@ \@minus 0.5\p@}
\setlength{\leftmargin}{3pc}
\setlength{\leftmargini}{\leftmargin}
\setlength{\leftmarginii}{2em}
\setlength{\leftmarginiii}{1.5em}
\setlength{\leftmarginiv}{1.0em}
\setlength{\leftmarginv}{0.5em}
\def\@listi  {\leftmargin\leftmargini}
\def\@listii {\leftmargin\leftmarginii
              \labelwidth\leftmarginii
              \advance\labelwidth-\labelsep
              \topsep  2\p@ \@plus 1\p@    \@minus 0.5\p@
              \parsep  1\p@ \@plus 0.5\p@ \@minus 0.5\p@
              \itemsep \parsep}

\providecommand{\maketitle}{}
\renewcommand{\maketitle}{%
  \par
  \begingroup
    \renewcommand{\thefootnote}{\fnsymbol{footnote}}
    \renewcommand{\@makefnmark}{\hbox to \z@{$^{\@thefnmark}$\hss}}
    \long\def\@makefntext##1{%
      \parindent 1em\noindent
      \hbox to 1.8em{\hss $\m@th ^{\@thefnmark}$}##1
    }
    \thispagestyle{empty}
    \@maketitle
    \@thanks
    \@notice
  \endgroup
  \let\maketitle\relax
  \let\thanks\relax
}

\newcommand{\@toptitlebar}{
  \hrule height 4\p@
  \vskip 0.25in
  \vskip -\parskip%
}
\newcommand{\@bottomtitlebar}{
  \vskip 0.29in
  \vskip -\parskip
  \hrule height 1\p@
  \vskip 0.09in%
}

\providecommand{\@maketitle}{}
\renewcommand{\@maketitle}{%
  \vbox{%
    \hsize\textwidth
    \linewidth\hsize
    \vskip 0.1in
    \@toptitlebar
    \centering
    {\LARGE\bf \@title\par}
    \@bottomtitlebar
    \if@submission
      \begin{tabular}[t]{c}\bf\rule{\z@}{24\p@}
        Anonymous Author(s) \\
        Affiliation \\
        Address \\
        \texttt{email} \\
      \end{tabular}%
    \else
      \def\And{%
        \end{tabular}\hfil\linebreak[0]\hfil%
        \begin{tabular}[t]{c}\bf\rule{\z@}{24\p@}\ignorespaces%
      }
      \def\AND{%
        \end{tabular}\hfil\linebreak[4]\hfil%
        \begin{tabular}[t]{c}\bf\rule{\z@}{24\p@}\ignorespaces%
      }
      \begin{tabular}[t]{c}\bf\rule{\z@}{24\p@}\@author\end{tabular}%
    \fi
    \vskip 0.3in \@minus 0.1in
  }
}

\newcommand{\ftype@noticebox}{8}
\newcommand{\@notice}{%
  \enlargethispage{2\baselineskip}%
  \@float{noticebox}[b]%
    \footnotesize\@noticestring%
  \end@float%
}

\renewenvironment{abstract}%
{%
  \vskip 0.075in%
  \centerline%
  {\large\bf Abstract}%
  \vspace{0.5ex}%
  \begin{quote}%
}
{
  \par%
  \end{quote}%
  \vskip 1ex%
}

\newcommand{\answerYes}[1][]{\textcolor{blue}{[Yes] #1}}
\newcommand{\answerNo}[1][]{\textcolor{orange}{[No] #1}}
\newcommand{\answerNA}[1][]{\textcolor{gray}{[NA] #1}}
\newcommand{\answerTODO}[1][]{\textcolor{red}{\bf [TODO]}}

\if@preprint
  \newcommand{\@noticestring}{%
    Preprint. Under review.%
  }
\else
  \if@neuripsfinal
    \newcommand{\@noticestring}{%
      \@neuripsordinal\/ Conference on Neural Information Processing Systems
      (NeurIPS \@neuripsyear).%
    }
  \else
    \newcommand{\@noticestring}{%
      Submitted to \@neuripsordinal\/ Conference on Neural Information
      Processing Systems (NeurIPS \@neuripsyear). Do not distribute.%
    }
  \fi
\fi
\endinput
\end{filecontents*}

\documentclass{article}

\PassOptionsToPackage{numbers,sort&compress}{natbib}
\usepackage[preprint]{neurips_2026}

\usepackage[utf8]{inputenc}
\usepackage[T1]{fontenc}
\usepackage{mathptmx}
\usepackage{amsmath,amssymb,amsthm,mathtools}
\usepackage{graphicx,booktabs,tabularx,array,caption,subcaption,xcolor,microtype}
\usepackage{hyperref}
\usepackage{url}

\definecolor{primary}{RGB}{18,48,120}
\definecolor{accent}{RGB}{0,110,90}
\definecolor{warn}{RGB}{175,65,25}
\definecolor{boxbg}{RGB}{244,247,252}
\definecolor{greenbg}{RGB}{240,250,246}
\definecolor{warnbg}{RGB}{253,245,240}
\definecolor{algobg}{RGB}{250,251,253}
\definecolor{bordergray}{RGB}{185,195,210}

\hypersetup{colorlinks=true, linkcolor=primary, citecolor=primary, urlcolor=primary}

\setlength{\parskip}{3.5pt plus 1pt minus 1pt}
\setlength{\parindent}{0pt}

\newtheorem{definition}{Definition}
\newtheorem{theorem}{Theorem}
\newtheorem{proposition}{Proposition}
\newtheorem{lemma}{Lemma}

\newcounter{algocf}
\renewcommand{\thealgocf}{\arabic{algocf}}

\newcommand{\E}{\mathbb{E}}
\newcommand{\R}{\mathbb{R}}
\newcommand{\IS}{\mathit{IS}}
\newcommand{\btheta}{\boldsymbol{\theta}}

\newcommand{\diagbox}[4]{%
  \fcolorbox{#1}{#2}{%
    \parbox[c]{#3}{\centering\footnotesize #4}%
  }%
}

\title{Interactive POMCP: Monte Carlo Tree Search for Multi-Agent Planning across Arbitrary Reasoning Levels}

\author{%
  Huy Truong \\
  Department of Computer Science \\
  \texttt{thuyng2@uic.edu} \\
}

\begin{document}

\maketitle

\begin{abstract}
Interactive Partially Observable Markov Decision Processes (I-POMDPs) provide a general Bayesian framework for sequential multi-agent decision-making under partial observability, in which an agent recursively maintains beliefs over both physical states and the intentional models of other agents across finitely nested reasoning levels ($L_0, L_1, \dots, L_k$). However, computing optimal policies in I-POMDPs suffers from the compounded curses of dimensionality and history because the interactive posterior's support grows exponentially with both the planning horizon and the nesting depth. Offline approximation techniques must either restrict the space of candidate opponent models a priori or evaluate branching reachability trees exponentially. In this paper, we introduce Interactive Partially Observable Monte Carlo Planning (I-POMCP), an online Monte Carlo tree search algorithm for finitely nested I-POMDPs across arbitrary reasoning levels. I-POMCP combines a canonical support-bounded interactive belief filter ($N=64$ atoms with low-variance systematic resampling) with an alternating history-action lookahead tree that maintains recursive lower and upper action-value bounds $[lo_r, hi_r]$ over unsampled observation branches. Empirical evaluations on the competitive and coordination Multi-Agent Tiger domains ($L_0$ to $L_4$) and the Multi-Agent Unmanned Aerial Vehicle (UAV) Reconnaissance benchmark ($L_1$ to $L_3$, horizons $H=2$ to $5$) demonstrate that I-POMCP converges to zero decision regret relative to exact dynamic programming solutions across all evaluated horizons, maintains $100\%$ ($1{,}016/1{,}016$) empirical interval containment, and outperforms particle-based Reachability Tree Sampling baselines across multi-step interactive episodes.
\end{abstract}

\section{Introduction}
\label{sec:intro}

Sequential decision-making in multi-agent environments requires an autonomous agent to reason not only about uncertain physical states, as in single-agent Partially Observable Markov Decision Processes (POMDPs) \citep{kaelbling1998planning,silver2010monte}, but also about the private beliefs, preferences, and strategic reasoning depths of other interacting agents. Interactive POMDPs (I-POMDPs), introduced by \citet{gmytrasiewicz2005framework}, formalize this multi-agent epistemic setting from an individual decision-maker's subjective perspective by augmenting the physical state space $S$ with finitely nested intentional models of other agents. Under a strategy level hierarchy $L_0, L_1, \dots, L_k$, a level-$l$ subject agent models its counterpart as a level-$(l-1)$ decision-maker and updates a hierarchical probability distribution over interactive states after every action and observation transition.

While the finitely nested I-POMDP formulation avoids the circularity of infinite belief hierarchies and permits non-equilibrium strategic reasoning, exact solution methods face a doubly exponential computational bottleneck. Even when the physical state space $S$, action spaces $A_i, A_j$, and observation spaces $\Omega_i, \Omega_j$ are finite, each single-step transition branches over all possible counterpart actions $a_j \in A_j$ and counterpart observations $o_j \in \Omega_j$. Consequently, the number of distinct level-$(l-1)$ counterpart belief distributions in the support of the level-$l$ posterior multiplies by up to $|A_j| \cdot |\Omega_j|$ at each timestep. Offline approximation methods such as Interactive Particle Filtering (IPF) with Reachability Tree Sampling (RTS) \citep{doshi2009monte} and behavioral equivalence compression \citep{zeng2012exploiting} alleviate offline intractability, yet they remain vulnerable to particle impoverishment when unexpected observations occur during execution.

In single-agent settings, Partially Observable Monte Carlo Planning (POMCP) \citep{silver2010monte} demonstrated that online Monte Carlo Tree Search (MCTS) \citep{kocsis2006bandit} can scale POMDP planning by focusing computation on the current belief state reached during execution. Extending online MCTS to finitely nested I-POMDPs across arbitrary reasoning levels $l \ge 1$, however, presents three fundamental algorithmic challenges. First, unweighted particle filters degenerate rapidly when propagated through nested opponent belief updates, because each particle must carry an entire recursive hierarchy of lower-level opponent beliefs. Second, evaluating exact opponent action distributions $\pi_j(\cdot \mid m_j)$ inside an online search tree requires recursively solving lower-level I-POMDPs at every encountered interactive state. Third, in domains with large observation branching factors, finite simulation budgets leave many observation branches unexplored, which can bias empirical Bellman backups unless unsampled probability mass is explicitly bounded.

\begin{figure}[t]
  \centering
  \fcolorbox{bordergray}{algobg}{%
    \begin{minipage}{0.97\textwidth}
      \vspace{4pt}
      \centering
      {\footnotesize\textbf{\color{primary}A. Finitely Nested Interactive Belief Hierarchy ($L_0 \to L_1 \to \dots \to L_l$)}}\\[4pt]
      \begin{tabular}{@{}c@{\hspace{5pt}}c@{\hspace{5pt}}c@{\hspace{5pt}}c@{\hspace{5pt}}c@{}}
        \diagbox{primary!60}{boxbg}{0.20\textwidth}{Level $L_0$ Counterpart\\$\IS_j^0 = S$, Uniform Policy\\$\pi_j^0(a_j) = 1/|A_j|$} &
        {\large $\xrightarrow{\text{models}}$} &
        \diagbox{primary!60}{boxbg}{0.22\textwidth}{Level $L_1$ Interactive State\\$\IS_i^1 = S \times M_j^0$\\$m_j^0 = (\btheta_j, \pi_j^0)$} &
        {\large $\xrightarrow{\text{recurses}}$} &
        \diagbox{primary}{boxbg}{0.23\textwidth}{Level $L_l$ Subject Agent\\$\IS_i^l = S \times M_j^{l-1}$\\$b_i^l \in \Delta(S \times (\btheta_j, b_j^{l-1}))$}
      \end{tabular}

      \vspace{6pt}
      {\large $\Downarrow$ \footnotesize\textit{Recursive Bayesian Belief Filtering \& Model Equivalence Merging}}
      \vspace{4pt}

      {\footnotesize\textbf{\color{accent}B. Interactive Belief Filtering Representations}}\\[4pt]
      \begin{tabular}{@{}c@{\hspace{6pt}}c@{\hspace{6pt}}c@{}}
        \diagbox{accent!70}{greenbg}{0.23\textwidth}{1. Exact Finite Filter\\Full posterior $\mathrm{SE}_\theta(b_i,a_i,o_i)$\\over canonical support;\\Exact reference oracle} &
        \diagbox{accent}{greenbg}{0.28\textwidth}{2. Compressed Filter (I-POMCP)\\Exact up to $N=64$ canonical atoms;\\Systematic resampling when\\posterior support $|\mathrm{supp}| > N$} &
        \diagbox{accent!70}{greenbg}{0.22\textwidth}{3. Interactive Particle Filter\\$N$-particle Monte Carlo\\propagation \citep{doshi2009monte};\\Paired with RTS baseline}
      \end{tabular}

      \vspace{6pt}
      {\large $\Downarrow$ \footnotesize\textit{Online Lookahead Search \& Certified Interval Backups}}
      \vspace{4pt}

      {\footnotesize\textbf{\color{warn}C. Online I-POMCP Planning \& Exact Reference Evaluation}}\\[4pt]
      \begin{tabular}{@{}c@{\hspace{6pt}}c@{\hspace{6pt}}c@{}}
        \diagbox{warn!80}{warnbg}{0.22\textwidth}{Recursive Solver Bank\\Memoized level-$(l-1)$\\policy queries $\pi_j(a_j \mid m_j)$} &
        \diagbox{warn!80}{warnbg}{0.28\textwidth}{Certified Value Intervals\\Recursive bounds $[\underline{Q}(h,a), \overline{Q}(h,a)]$\\using $[lo_r, hi_r]$ on unsampled $o_i$} &
        \diagbox{warn!80}{warnbg}{0.23\textwidth}{Exact Reference Comparison\\$\alpha$-vector DP ($L_1$) \& exact\\expectation trees ($L_1$ to $L_3$)}
      \end{tabular}
      \vspace{4pt}
    \end{minipage}%
  }
  \caption{Overview of the Interactive POMCP (I-POMCP) framework. (A) Finitely nested interactive state hierarchy from level $L_0$ to arbitrary reasoning level $L_l$. (B) Belief filtering representations comparing exact finite filtering, canonical support-bounded compressed filtering ($N=64$), and Interactive Particle Filtering. (C) Online MCTS lookahead with memoized lower-level opponent policy queries and recursive certified value bounds.}
  \label{fig:ipomcp_arch}
\end{figure}

In this paper, we present Interactive POMCP (I-POMCP), an online Monte Carlo planning framework for finitely nested I-POMDPs across arbitrary reasoning levels (Figure~\ref{fig:ipomcp_arch}). Rather than relying on unweighted particle propagation, I-POMCP maintains a compressed canonical belief distribution that computes exact recursive Bayesian posteriors whenever the merged support size does not exceed a capacity threshold $N$ (default $N=64$), and applies low-variance systematic resampling only when posterior branching exceeds $N$. During online tree search, I-POMCP combines Upper Confidence Bound for Trees (UCT) exploration \citep{kocsis2006bandit} with exact recursive value bounds $[lo_r, hi_r]$ that account for unsampled observation branches, enabling certified action selection and direct verification against exact dynamic programming solvers.

The remainder of this paper is organized as follows. Section~\ref{sec:background} reviews the finitely nested I-POMDP formulation, the recursive Bayesian belief update operator, and the baseline policy classes considered in this study. Section~\ref{sec:method} introduces the compressed interactive belief filter (Algorithm~\ref{alg:compressed_filter}) and the online I-POMCP tree search algorithm with recursive value bounds (Figure~\ref{fig:mcts_tree} and Algorithm~\ref{alg:ipomcp_search}). Section~\ref{sec:experiments} presents our empirical evaluation on Multi-Agent Tiger ($L_0$ to $L_4$) and Multi-Agent UAV Reconnaissance ($L_1$ to $L_3$, horizons $H=2$ to $5$), comparing I-POMCP against exact dynamic programming reference solvers, Interactive Particle Filtering with Reachability Tree Sampling, exact-belief heuristics, reactive policies, and random walks. Section~\ref{sec:conclusion} concludes the paper.

\section{Background and Finitely Nested Interactive POMDPs}
\label{sec:background}

We consider a sequential interaction between a subject agent $i$ and an interacting counterpart agent $j$. Following \citet{gmytrasiewicz2005framework}, a finitely nested I-POMDP of strategy level $l \ge 1$ for agent $i$ is defined by the tuple:
\begin{equation}
  \mathrm{I\text{-}POMDP}_i^l = \left\langle \IS_i^l, \, A, \, T_i, \, \Omega_i, \, O_i, \, R_i, \, \gamma, \, H \right\rangle,
  \label{eq:ipomdp_tuple}
\end{equation}
where $\IS_i^l$ is the space of level-$l$ interactive states, $A = A_i \times A_j$ is the finite set of joint physical actions, $\Omega_i$ is the finite set of observations available to agent $i$, $R_i: S \times A_i \times A_j \to [r_{\min}, r_{\max}]$ is a bounded immediate reward function, $\gamma \in (0, 1]$ is the discount factor, and $H \in \mathbb{N}$ is the planning horizon.

\subsection{Interactive State Space Hierarchy and Recursive Belief Update}

To reason strategically about counterpart $j$ without assuming common knowledge of a pre-computed equilibrium, agent $i$ maintains a finitely nested belief hierarchy grounded at level $0$.

\begin{definition}[Finitely Nested Interactive State Space]
  For a subject agent $i$ and counterpart agent $j$, the level-$0$ interactive state space coincides with the finite physical state space, $\IS_i^0 = S$. For any reasoning level $l \ge 1$, the level-$l$ interactive state space $\IS_i^l$ is defined inductively as:
  \begin{equation}
    \IS_i^l = S \times M_j^{l-1}, \qquad M_j^{l-1} = \left\{ m_j = (\btheta_j, b_j^{l-1}) : b_j^{l-1} \in \Delta(\IS_j^{l-1}) \right\},
    \label{eq:is_hierarchy}
  \end{equation}
  where $\btheta_j = (T_j, A_j, \Omega_j, O_j, R_j, \gamma, H)$ denotes the intentional frame of agent $j$ and $b_j^{l-1}$ denotes the level-$(l-1)$ probability distribution (belief) of agent $j$ over $\IS_j^{l-1}$. A level-$l$ belief of agent $i$ is a probability distribution $b_i^l \in \Delta(\IS_i^l)$.
  \label{def:is_space}
\end{definition}

Under the standard physical state sufficiency and model non-manipulability assumptions of \citet{gmytrasiewicz2005framework}, the physical state transitions according to the Markov kernel $T(s' \mid s, a_i, a_j)$, and agents $i$ and $j$ receive observations drawn independently from $O_i(o_i \mid s', a_i, a_j)$ and $O_j(o_j \mid s', a_i, a_j)$, respectively. When agent $i$ at level $l > 0$ executes action $a_i \in A_i$ and receives observation $o_i \in \Omega_i$, its updated posterior belief $b_i'(s', m_j') = \mathrm{SE}_\theta(b_i, a_i, o_i)(s', m_j')$ marginalizes over unobserved counterpart actions $a_j \in A_j$ and unobserved counterpart observations $o_j \in \Omega_j$, recursively invoking the counterpart's level-$(l-1)$ belief update operator $U_j(m_j, a_j, o_j) = (\btheta_j, \mathrm{SE}_{\theta_j}(b_j^{l-1}, a_j, o_j))$:
\begin{equation}
  \begin{aligned}
    b_i'(s', m_j') = \frac{1}{P(o_i \mid b_i, a_i)} \sum_{s, m_j} b_i(s, m_j) \sum_{a_j \in A_j} \pi_j(a_j \mid m_j) \, T(s' \mid s, a_i, a_j) \, O_i(o_i \mid s', a_i, a_j) \\
    \times \sum_{o_j \in \Omega_j} O_j(o_j \mid s', a_i, a_j) \, \mathbb{I}\!\left[m_j' = U_j(m_j, a_j, o_j)\right],
  \end{aligned}
  \label{eq:exact_se}
\end{equation}
where $P(o_i \mid b_i, a_i)$ is the observation normalizing constant and $\pi_j(a_j \mid m_j)$ is the level-$(l-1)$ optimal action distribution obtained by solving model $m_j$ over the remaining horizon.

\subsection{Optimal Value Functions and Baseline Decision Policies}

Given a level-$l$ belief $b_i \in \Delta(\IS_i^l)$ and remaining horizon $r \ge 1$, the exact finite-horizon optimal value function $V_r^*(b_i) = \max_{a_i \in A_i} Q_r^*(b_i, a_i)$ satisfies the interactive Bellman optimality equation \citep{gmytrasiewicz2005framework}:
\begin{equation}
  Q_r^*(b_i, a_i) = \sum_{s, m_j} b_i(s, m_j) \sum_{a_j \in A_j} \pi_j(a_j \mid m_j) R_i(s, a_i, a_j) + \gamma \sum_{o_i \in \Omega_i} P(o_i \mid b_i, a_i) V_{r-1}^*\!\left(\mathrm{SE}_\theta(b_i, a_i, o_i)\right),
  \label{eq:bellman}
\end{equation}
with base condition $V_0^*(b_i) = 0$. When counterpart $j$ has multiple tied optimal actions, $\pi_j(a_j \mid m_j)$ assigns uniform probability across the maximizers of $Q_{r}^{\theta_j}(b_j^{l-1}, a_j)$. To evaluate online planning across reasoning depths, we formally define the baseline and reference policy classes compared in our study.

\begin{definition}[Baseline and Reference Decision Policies]
  \sloppy Let $h_t = (a_i^0, o_i^1, \dots, a_i^{t-1}, o_i^t)$ denote the action-observation history of agent $i$ at timestep $t$, and let $b_i^t \in \Delta(\IS_i^l)$ denote its posterior belief. We distinguish five canonical policy classes:
  \begin{enumerate}
    \item A random walk policy (or level-$0$ uniform policy, $\pi^0$) holds no interactive belief and selects actions uniformly at random at every timestep, $\pi^0(a_i \mid h_t) = 1/|A_i|$ for all $a_i \in A_i$. This policy serves both as the base of the recursive hierarchy ($L_0$) and as a zero-information baseline.
    \item A reactive policy ($\pi^{\mathrm{react}}: \Omega_i \to \Delta(A_i)$) is a memoryless mapping that conditions action selection solely on the most recent immediate observation $o_i^t \in \Omega_i$, discarding prior history $h_{t-1}$ and performing no sequential Bayesian filtering.
    \item An exact-belief heuristic policy ($\pi^{\mathrm{heur}}: \Delta(\IS_i^l) \to A_i$) maintains the exact recursive Bayesian belief $b_i^t$ via Equation~\eqref{eq:exact_se}, but maps $b_i^t$ to an action using a fixed domain-specific threshold or myopic rule rather than multi-step lookahead search (for example, listening in Multi-Agent Tiger until the marginal posterior probability of the tiger's location exceeds a fixed confidence threshold, or advancing greedily toward the highest-probability target cell in UAV Reconnaissance).
    \item An Interactive Particle Filter with Reachability Tree Sampling policy ($\mathrm{IPF+RTS}$) \citep{doshi2009monte} approximates the posterior $b_i^t$ using a set of $N$ Monte Carlo particles and computes finite-horizon action values by forward-sampling reachable particle belief trees of depth $H$.
    \item An exact reference solver computes the true Bayes-optimal action value $Q_H^*(b_i^t, a_i)$ from Equation~\eqref{eq:bellman}, either via exact finite-horizon $\alpha$-vector value iteration at level $L_1$ (where the fixed level-$0$ counterpart policy $\pi_j^0$ reduces the environment to a standard POMDP) or via exact recursive expectation tree evaluation using the exact finite filter at levels $L_1$ to $L_3$.
  \end{enumerate}
  \label{def:baselines}
\end{definition}

\begin{figure}[t]
  \refstepcounter{algocf}\label{alg:compressed_filter}
  \centering
  \fcolorbox{primary!50}{algobg}{%
    \begin{minipage}{0.97\textwidth}
      \vspace{3pt}
      {\small\textbf{\color{primary}Algorithm \thealgocf: \texttt{CompressInteractiveBelief}, Canonical Support-Bounded Recursive Filter ($L_l$)}}\\[2pt]
      \rule{\textwidth}{0.6pt}\\[-2pt]
      \footnotesize
      Require: Level-$l$ interactive belief $b_i^l = \{( (s_k, m_{j,k}), w_k )\}_{k=1}^K$, subject action $a_i \in A_i$, observation $o_i \in \Omega_i$, atom capacity $N$ (default $N=64$).\\
      Ensure: Updated canonical posterior belief $\hat{b}_i^{\prime l}$ with $|\mathrm{supp}(\hat{b}_i^{\prime l})| \le N$.\\[2pt]
      \begin{tabular}{@{}r@{\hspace{6pt}}p{0.90\textwidth}@{}}
        1: & Initialize canonical hash map $\mathcal{M} \leftarrow \emptyset$ mapping canonical interactive state keys $\mathrm{key}(s', m_j') \mapsto \text{unnormalized weight}$. \\
        2: & \textbf{for each} support atom $((s, m_j), w) \in b_i^l$ \textbf{with} $w > 0$ \textbf{do} \\
        3: & \hspace{1.2em} Query recursive level-$(l-1)$ counterpart policy distribution $\pi_j(\cdot \mid m_j)$ from memoized solver bank. \\
        4: & \hspace{1.2em} \textbf{for each} counterpart action $a_j \in A_j$ \textbf{with} $\pi_j(a_j \mid m_j) > 0$ \textbf{and} observation $o_j \in \Omega_j$ \textbf{do} \\
        5: & \hspace{2.4em} Compute updated counterpart model $m_j' \leftarrow U_j(m_j, a_j, o_j)$ recursively via $\mathrm{CompressInteractiveBelief}$ at level $l-1$. \\
        6: & \hspace{2.4em} \textbf{for each} next physical state $s' \in S$ \textbf{with} $T(s' \mid s, a_i, a_j) > 0$ \textbf{do} \\
        7: & \hspace{3.6em} $\Delta w \leftarrow w \cdot \pi_j(a_j \mid m_j) \cdot T(s' \mid s, a_i, a_j) \cdot O_i(o_i \mid s', a_i, a_j) \cdot O_j(o_j \mid s', a_i, a_j)$. \\
        8: & \hspace{3.6em} \textbf{if} $\Delta w > 0$ \textbf{then} accumulate $\mathcal{M}[\mathrm{key}(s', m_j')] \leftarrow \mathcal{M}[\mathrm{key}(s', m_j')] + \Delta w$. \\
        9: & \hspace{2.4em} \textbf{end for} \\
        10: & \hspace{1.2em} \textbf{end for} \\
        11: & \textbf{end for} \\
        12: & Normalize weights in $\mathcal{M}$ by $Z = \sum_{(s', m_j') \in \mathcal{M}} \mathcal{M}[\mathrm{key}(s', m_j')]$ to form exact merged posterior $\tilde{b}_i'$. \\
        13: & \textbf{if} $|\mathrm{supp}(\tilde{b}_i')| \le N$ \textbf{then} \\
        14: & \hspace{1.2em} \textbf{return} $\hat{b}_i^{\prime l} \leftarrow \tilde{b}_i'$ \hfill $\triangleright$ \textit{Exact posterior retention (Proposition~\ref{prop:compressed_exact})} \\
        15: & \textbf{else} \\
        16: & \hspace{1.2em} Sort atoms of $\tilde{b}_i'$ by canonical key and construct cumulative distribution $F_1, \dots, F_{|\mathcal{M}|}$. \\
        17: & \hspace{1.2em} Sample single offset $u_0 \sim \mathrm{Uniform}[0, 1/N)$ and set systematic grid $u_m = u_0 + \frac{m-1}{N}$ for $m = 1, \dots, N$. \\
        18: & \hspace{1.2em} Aggregate selected atoms into canonical weighted set $\hat{b}_i^{\prime l}$ of size $\le N$ and \textbf{return} $\hat{b}_i^{\prime l}$. \\
        19: & \textbf{end if}
      \end{tabular}
      \vspace{3pt}
    \end{minipage}%
  }
\end{figure}

\section{Interactive Monte Carlo Planning (I-POMCP)}
\label{sec:method}

Directly evaluating Equation~\eqref{eq:exact_se} inside an online lookahead tree of horizon $H$ causes the support size $|\mathrm{supp}(b_i)|$ to grow exponentially with depth, whereas standard particle filtering \citep{doshi2009monte} introduces unnecessary Monte Carlo variance even when the reachable interactive state space is small. I-POMCP addresses both issues by combining a canonical support-bounded belief filter (Algorithm~\ref{alg:compressed_filter}) with an alternating history-action lookahead tree equipped with recursive value bounds (Figure~\ref{fig:mcts_tree} and Algorithm~\ref{alg:ipomcp_search}).

\subsection{Support-Bounded Compressed Filtering via Canonical Merging and Systematic Resampling}

A key structural observation in finitely nested I-POMDPs is that distinct joint physical transitions $(s, a_i, a_j) \to s'$ often induce identical posterior counterpart beliefs $b_j^{\prime l-1}$ whenever the counterpart's observation $o_j$ depends only on $(a_i, a_j)$ or $s'$. As shown in Algorithm~\ref{alg:compressed_filter}, our compressed filter first enumerates the exact one-step transitions from the current support atoms and merges all transitions that map to the same canonical interactive state $(s', m_j')$. If the resulting merged posterior $\tilde{b}_i'$ contains at most $N$ distinct atoms (default $N=64$), the filter retains $\tilde{b}_i'$ exactly. Only when $|\mathrm{supp}(\tilde{b}_i')| > N$ does the filter compress the support to at most $N$ atoms via low-variance systematic resampling.

\begin{proposition}[Exactness under Bounded Canonical Support]
  Let $b_0$ be an initial level-$l$ interactive belief with $|\mathrm{supp}(b_0)| \le N$. For any action-observation history $h_t = (a_i^0, o_i^1, \dots, a_i^{t-1}, o_i^t)$ along which the canonical support after merging equivalent counterpart models satisfies $|\mathrm{supp}(\mathrm{SE}^t(b_0, h_t))| \le N$ at each step $1 \le \tau \le t$, the compressed filter output $\hat{b}_t^{(N)}$ coincides exactly with the true Bayesian posterior, satisfying $\|\hat{b}_t^{(N)} - \mathrm{SE}^t(b_0, h_t)\|_1 = 0$.
  \label{prop:compressed_exact}
\end{proposition}

\begin{figure}[t]
  \centering
  \fcolorbox{bordergray}{algobg}{%
    \begin{minipage}{0.97\textwidth}
      \vspace{4pt}
      \centering
      {\footnotesize\textbf{\color{primary}Alternating History-Action Lookahead Tree with Certified Unsampled Branch Bounds}}\\[4pt]
      \diagbox{primary}{boxbg}{0.52\textwidth}{Root History Node $h$ (Remaining Horizon $r$)\\Canonical Compressed Belief $\hat{b}_h^l = \{((s_k, m_{j,k}), w_k)\}_{k=1}^{\le 64}$}\\[3pt]
      {\small $\swarrow$ \hspace{8.5em} $\downarrow$ \hspace{8.5em} $\searrow$}\\[2pt]
      \begin{tabular}{@{}c@{\hspace{8pt}}c@{\hspace{8pt}}c@{}}
        \diagbox{primary!70}{boxbg}{0.22\textwidth}{Action Branch $(h, a_1)$\\Visits $N(h,a_1) = 42$\\Mean Reward $\bar{R}(h,a_1)$} &
        \diagbox{primary}{boxbg}{0.28\textwidth}{Selected Action $(h, a^*)$\\Visits $N(h,a^*) = 128$\\$Q(h,a^*) = \bar{R} + \gamma \sum_o \frac{C(h,a^*,o)}{N(h,a^*)} V(ha^*o)$} &
        \diagbox{primary!70}{boxbg}{0.22\textwidth}{Action Branch $(h, a_3)$\\Visits $N(h,a_3) = 18$\\Pruned if $\overline{Q}(h,a_3) < \underline{Q}(h,a^*)$}
      \end{tabular}\\[3pt]
      {\small $\swarrow$ \hspace{5.5em} $\searrow$}\\[2pt]
      \begin{tabular}{@{}c@{\hspace{12pt}}c@{}}
        \diagbox{accent}{greenbg}{0.42\textwidth}{Sampled Observation Branch $ha^*o_1$ ($C(h,a^*,o_1) > 0$)\\Updated belief $\hat{b}_{ha^*o_1} \leftarrow \mathrm{CompressInteractiveBelief}(\hat{b}_h, a^*, o_1)$\\Contributes weight $P(o_1 \mid \hat{b}_h, a^*) \cdot [\underline{V}(ha^*o_1), \overline{V}(ha^*o_1)]$} &
        \diagbox{warn}{warnbg}{0.42\textwidth}{Unsampled Observation Branch $ha^*o_2$ ($C(h,a^*,o_2) = 0$)\\No child rollout yet; bounded conservatively via reward range:\\Contributes prior mass $P(o_2 \mid \hat{b}_h, a^*) \cdot [lo_{r-1}, \, hi_{r-1}]$}
      \end{tabular}
      \vspace{4pt}
    \end{minipage}%
  }
  \caption{Online lookahead tree structure in I-POMCP. Sampled observation children back up recursive value estimates and certified sub-intervals, while unsampled observation branches ($C(h,a,o)=0$) contribute worst- and best-case horizon bounds $[lo_{r-1}, hi_{r-1}]$ scaled by their prior probability $P(o \mid \hat{b}_h, a)$.}
  \label{fig:mcts_tree}
\end{figure}

\subsection{Online Alternating Tree Search and Recursive Certified Value Bounds}

At each decision step, I-POMCP constructs an alternating history-action lookahead tree rooted at the current belief $b_0^l$ and remaining horizon $H$ (Figure~\ref{fig:mcts_tree} and Algorithm~\ref{alg:ipomcp_search}). Each history node $h$ stores its canonical compressed belief $\hat{b}_h^l$, visit count $N(h)$, action visit counts $N(h, a)$, observation branch counts $C(h, a, o)$ satisfying $\sum_{o \in \Omega_i} C(h, a, o) = N(h, a)$, and running mean immediate reward $\bar{R}(h, a)$. During each simulation, actions are selected via the UCT rule \citep{kocsis2006bandit,silver2010monte}, and action values are backed up via the empirical Bellman operator:
\begin{equation}
  Q(h, a) = \bar{R}(h, a) + \gamma \sum_{o \in \Omega_i} \frac{C(h, a, o)}{N(h, a)} V(hao), \qquad V(h) = \max_{a \in A_i} Q(h, a).
  \label{eq:ipomcp_backup}
\end{equation}

In finite-sample search over wide observation spaces, some observation branches $o \in \Omega_i$ may remain unsampled ($C(h, a, o) = 0$) or only partially explored. To bound the potential value contribution of unexplored branches and provide a rigorous stopping certificate, we maintain recursive horizon value intervals.

\begin{definition}[Recursive Certified Value Bounds and Evaluation Metrics]
  \sloppy Let $[r_{\min}, r_{\max}]$ denote the immediate reward bounds of $R_i$. For remaining rollout horizon $r \ge 1$, the worst-case and best-case cumulative return bounds $[lo_r, hi_r]$ are defined recursively by:
  \begin{equation}
    lo_r = r_{\min} + \gamma \min(0, lo_{r-1}), \qquad hi_r = r_{\max} + \gamma \max(0, hi_{r-1}), \qquad lo_0 = hi_0 = 0.
    \label{eq:certified_bounds}
  \end{equation}
  At any tree node $h$ with remaining horizon $r$ and canonical belief $\hat{b}_h$, let $\underline{V}(hao) = \max_{a'} \underline{Q}(hao, a')$ and $\overline{V}(hao) = \max_{a'} \overline{Q}(hao, a')$ for sampled children ($C(h,a,o) > 0$), and set $[\underline{V}(hao), \overline{V}(hao)] = [lo_{r-1}, hi_{r-1}]$ for unsampled branches ($C(h,a,o)=0$). Weighting by the belief-induced immediate reward and observation probabilities under $\hat{b}_h$ yields the recursive certified action-value interval:
  \begin{equation}
    \begin{aligned}
      \underline{Q}(h, a) &= \E_{\hat{b}_h}[R_i(s, a, a_j)] + \gamma \sum_{o \in \Omega_i} P(o \mid \hat{b}_h, a) \underline{V}(hao), \\
      \overline{Q}(h, a) &= \E_{\hat{b}_h}[R_i(s, a, a_j)] + \gamma \sum_{o \in \Omega_i} P(o \mid \hat{b}_h, a) \overline{V}(hao).
    \end{aligned}
    \label{eq:interval_backup}
  \end{equation}
  Given exact reference action values $Q_H^*(b_0, a)$ and optimal action $a^* = \arg\max_a Q_H^*(b_0, a)$, we evaluate three metrics for the action $\hat{a} = \arg\max_a Q(h_0, a)$ selected by I-POMCP:
  \begin{enumerate}
    \item Decision Regret (Planner Loss), defined as $\mathcal{L}_{\mathrm{plan}}(b_0, \hat{a}) = Q_H^*(b_0, a^*) - Q_H^*(b_0, \hat{a}) \ge 0$, which measures the true expected return suboptimality of the chosen action $\hat{a}$ ($\mathcal{L}_{\mathrm{plan}} = 0.00$ iff $\hat{a}$ is Bayes-optimal).
    \item Optimal Action-Value Error (\texttt{optimal\_q\_error}), defined as $|Q(h_0, a^*) - Q_H^*(b_0, a^*)|$, which measures value estimation accuracy at the true optimal action.
    \item Certified Interval Containment, defined as the indicator $\mathbb{I}[Q_H^*(b_0, a) \in [\underline{Q}(h_0, a), \overline{Q}(h_0, a)] \;\forall a \in A_i]$.
  \end{enumerate}
  \label{def:metrics}
\end{definition}

\begin{figure}[t]
  \refstepcounter{algocf}\label{alg:ipomcp_search}
  \centering
  \fcolorbox{primary!50}{algobg}{%
    \begin{minipage}{0.97\textwidth}
      \vspace{3pt}
      {\small\textbf{\color{primary}Algorithm \thealgocf: \texttt{IPOMCP-Search}, Arbitrary-Level Online MCTS with Certified Horizon Bounds}}\\[2pt]
      \rule{\textwidth}{0.6pt}\\[-2pt]
      \footnotesize
      Require: Root level-$l$ belief $b_0^l$, planning horizon $H$, simulation budget $M_{\text{sim}}$, exploration constant $c_{\text{ucb}}$, atom capacity $N=64$.\\
      Ensure: Selected action $a^* \in A_i$ and root value estimate $V(h_0)$.\\[2pt]
      \begin{tabular}{@{}r@{\hspace{6pt}}p{0.90\textwidth}@{}}
        1: & Precompute horizon return bounds $[lo_r, hi_r]$ for $r = 0, 1, \dots, H$ via Eq.~\eqref{eq:certified_bounds}. \\
        2: & Initialize root history node $h_0 \leftarrow \emptyset$ with canonical belief $\hat{b}_{h_0} \leftarrow b_0^l$. \\
        3: & \textbf{for} simulation $m = 1, \dots, M_{\text{sim}}$ \textbf{do} \\
        4: & \hspace{1.2em} Sample interactive state $(s, m_j) \sim \hat{b}_{h_0}$ and invoke $\textsc{Simulate}(h_0, (s, m_j), H)$. \\
        5: & \hspace{1.2em} \textbf{if} $\min_{a \neq \hat{a}^*} \!\left(\underline{Q}(h_0, \hat{a}^*) - \overline{Q}(h_0, a)\right) > 0$ \textbf{then break} \hfill $\triangleright$ \textit{Early stopping on interval separation} \\
        6: & \textbf{end for} \\
        7: & \textbf{return} $a^* \leftarrow \arg\max_{a \in A_i} Q(h_0, a)$. \\[3pt]
        8: & \textbf{procedure} $\textsc{Simulate}(h, (s, m_j), r)$ \\
        9: & \hspace{1.2em} \textbf{if} $r = 0$ \textbf{then return} $0$. \\
        10: & \hspace{1.2em} Select subject action $a_i \leftarrow \arg\max_{a \in A_i} \!\left( Q(h, a) + c_{\text{ucb}} \sqrt{\frac{\ln(N(h)+1)}{N(h,a)+1}} \right)$. \\
        11: & \hspace{1.2em} Sample counterpart action $a_j \sim \pi_j(\cdot \mid m_j)$, next state $s' \sim T(\cdot \mid s, a_i, a_j)$, and observations $o_i \sim O_i$, $o_j \sim O_j$. \\
        12: & \hspace{1.2em} Compute immediate reward $R \leftarrow R_i(s, a_i, a_j)$; if child $ha_io_i$ is new, initialize $\hat{b}_{ha_io_i}$ via Algorithm~\ref{alg:compressed_filter}. \\
        13: & \hspace{1.2em} Recursively evaluate child node $G_{\text{child}} \leftarrow \textsc{Simulate}(ha_io_i, (s', U_j(m_j, a_j, o_j)), r-1)$. \\
        14: & \hspace{1.2em} Increment counts $N(h) \leftarrow N(h)+1$, $N(h,a_i) \leftarrow N(h,a_i)+1$, $C(h,a_i,o_i) \leftarrow C(h,a_i,o_i)+1$. \\
        15: & \hspace{1.2em} Update running mean reward $\bar{R}(h,a_i)$ with $R$, recompute $Q(h,a_i)$ via Eq.~\eqref{eq:ipomcp_backup}, and update $[\underline{Q}(h,a_i), \overline{Q}(h,a_i)]$. \\
        16: & \hspace{1.2em} \textbf{return} $R + \gamma V(ha_io_i)$. \\
        17: & \textbf{end procedure}
      \end{tabular}
      \vspace{3pt}
    \end{minipage}%
  }
\end{figure}

\begin{theorem}[Interval Containment and Asymptotic Convergence of I-POMCP]
  Suppose the initial level-$l$ belief $b_0^l$ and all reachable posteriors over horizon $H$ have canonical support size at most $N$, and suppose lower-level counterpart policies $\pi_j(\cdot \mid m_j)$ are evaluated exactly. Then: (1) for any simulation count $M_{\mathrm{sim}} \ge 1$, the certified interval $[\underline{Q}(h_0, a), \overline{Q}(h_0, a)]$ in Equation~\eqref{eq:interval_backup} satisfies $\underline{Q}(h_0, a) \le Q_H^*(b_0^l, a) \le \overline{Q}(h_0, a)$ for all $a \in A_i$; and (2) as $M_{\mathrm{sim}} \to \infty$, the I-POMCP action-value estimate $Q(h_0, a)$ converges in probability to $Q_H^*(b_0^l, a)$ for all $a \in A_i$, and $\mathcal{L}_{\mathrm{plan}}(b_0^l, \hat{a}) \xrightarrow{p} 0$.
  \label{thm:ipomcp_conv}
\end{theorem}

\section{Experimental Evaluation}
\label{sec:experiments}

We evaluate I-POMCP on two standard multi-agent benchmark domains from the I-POMDP literature \citep{gmytrasiewicz2005framework,doshi2009monte,zeng2012exploiting}: the Multi-Agent Tiger Problem and the Multi-Agent UAV Reconnaissance Problem (Table~\ref{tab:domains}).

\subsection{Benchmark Domains and Nested Initial Belief Configurations}

In the Multi-Agent Tiger domain ($|S|=2$, $|A_i|=3$, $|\Omega_i|=2$), two agents stand before two closed doors hiding a tiger ($T_L$ or $T_R$) and treasure. Each agent may listen ($\mathrm{L}$) or open the left or right door ($\mathrm{O}_L, \mathrm{O}_R$), receiving noisy acoustic observations regarding the tiger's position as well as noisy creak observations indicating whether the counterpart opened a door. We evaluate both the competitive (zero-sum / conflicting reward) and coordination (cooperative joint door-opening bonus) payoff structures across reasoning levels $L_0$ through $L_4$.

In the Multi-Agent UAV Reconnaissance domain \citep{doshi2009monte,zeng2012exploiting}, an autonomous reconnaissance UAV (agent $i$) navigates a grid environment ($|A_i|=5$ movement and hover actions, $|\Omega_i|=3$ sensor observations) to intercept or track an evading target agent $j$ under partial visibility. Because UAV Reconnaissance has $5$ actions and $3$ observations per step (compared to Tiger's $3$ actions and $2$ observations, plus terminal spotting transitions), the lookahead tree at horizon $H=5$ is approximately $15\times$ wider than in Tiger.

To evaluate planners systematically across varied initial epistemic conditions, we construct nested initial belief hierarchies from a canonical set of $4$ distinct initial physical state distributions $\mathcal{B}_0 = \{b_{[1]}, b_{[2]}, b_{[3]}, b_{[4]}\}$ per domain (Figure~\ref{fig:scaling_panel}A). A level-$l$ problem instance is specified by a tuple $(b^{(1)}, b^{(2)}, \dots, b^{(l)}) \in \mathcal{B}_0^l$ assigning the initial physical prior at each nested level $1, \dots, l$. This construction yields $4^1 = 4$ problem configurations at $L_1$, $4^2 = 16$ configurations at $L_2$, and $4^3 = 64$ configurations at $L_3$.

\begin{table}[t]
  \caption{Summary of Benchmark Domains, Evaluated Policy Classes, and Structural Properties. All baseline and reference policies ($\pi^0$, $\pi^{\mathrm{react}}$, $\pi^{\mathrm{heur}}$, $\mathrm{IPF+RTS}$, and exact reference solvers) are formally defined in Definition~\ref{def:baselines}.}
  \label{tab:domains}
  \centering
  \small
  \begin{tabularx}{\textwidth}{>{\raggedright\arraybackslash}p{0.21\textwidth} >{\raggedright\arraybackslash}p{0.19\textwidth} >{\raggedright\arraybackslash}p{0.21\textwidth} >{\raggedright\arraybackslash}X}
    \toprule
    \textbf{Domain} & \textbf{State / Action / Obs.} & \textbf{Levels \& Policies} & \textbf{Domain Structure \& Evaluation Scope} \\
    \midrule
    Multi-Agent Tiger & $|S|=2$, $|A_i|=3$, $|\Omega_i|=2$ & $L_0$ to $L_4$: $\pi^0$, $\pi^{\mathrm{react}}$, $\pi^{\mathrm{heur}}$, $\mathrm{IPF+RTS}$, $\mathrm{I\text{-}POMCP}$ & Competitive and coordination payoff matrices; noisy listening and door-creak observations; evaluated across $4^l$ nested initial belief hierarchies and multi-episode interactions. \\
    \addlinespace[3pt]
    UAV Reconnaissance & Grid state, $|A_i|=5$, $|\Omega_i|=3$ & $L_1$ to $L_3$: $\pi^0$, $\pi^{\mathrm{heur}}$, $\mathrm{IPF+RTS}$, $\mathrm{I\text{-}POMCP}$ & Reconnaissance game of \citet{doshi2009monte} and \citet{zeng2012exploiting}; $\sim 15\times$ wider lookahead tree at $H=5$; evaluated against exact recursive solvers up to $H=5$. \\
    \bottomrule
  \end{tabularx}
\end{table}

\begin{figure}[t]
  \centering
  \fcolorbox{bordergray}{algobg}{%
    \begin{minipage}{0.97\textwidth}
      \vspace{4pt}
      \begin{tabular}{@{}>{\raggedright\arraybackslash}p{0.46\textwidth}@{\hspace{10pt}}>{\raggedright\arraybackslash}p{0.49\textwidth}@{}}
        {\footnotesize\textbf{\color{primary}A. Nested Initial Belief Hierarchy Construction ($4^l$)}} &
        {\footnotesize\textbf{\color{accent}B. Strategic Depth Saturation on Coordination Tiger ($L_1$ to $L_4$)}}\\[4pt]
        \fcolorbox{primary!40}{boxbg}{%
          \begin{minipage}{0.43\textwidth}
            \raggedright\footnotesize
            Canonical Physical Priors: $\mathcal{B}_0 = \{b_{[1]}, b_{[2]}, b_{[3]}, b_{[4]}\}$\\[3pt]
            $\bullet$ Level $L_1$ ($4^1 = 4$ instances): Prior tuple $(b^{(1)}) \in \mathcal{B}_0$\\[2pt]
            $\bullet$ Level $L_2$ ($4^2 = 16$ instances): Tuple $(b^{(1)}, b^{(2)}) \in \mathcal{B}_0^2$\\[2pt]
            $\bullet$ Level $L_3$ ($4^3 = 64$ instances): Tuple $(b^{(1)}, b^{(2)}, b^{(3)}) \in \mathcal{B}_0^3$\\[3pt]
            Each instance defines a complete finitely nested initial belief $b_0^l \in \Delta(\IS_i^l)$ for exact reference evaluation.
          \end{minipage}%
        } &
        \fcolorbox{accent!40}{greenbg}{%
          \begin{minipage}{0.46\textwidth}
            \raggedright\footnotesize
            Action Agreement with $L_3$ Policy \& Compute Cost:\\[4pt]
            $L_1$ vs.\ $L_3$: \colorbox{primary!30}{\rule{18pt}{5pt}} Low agreement ($<0.1\,\text{CPU-h}$)\\[3pt]
            $L_2$ vs.\ $L_3$: \colorbox{primary!55}{\rule{45pt}{5pt}} Moderate ($<1.0\,\text{CPU-h}$)\\[3pt]
            $L_3$ (Ref.): \colorbox{primary!85}{\rule{78pt}{5pt}} $100\%$ ($40/40$ trials)\\[3pt]
            $L_4$ vs.\ $L_3$: \colorbox{warn!80}{\rule{68pt}{5pt}} $85.0\text{--}87.5\%$ ($34\text{--}35/40$, $69\,\text{CPU-h}$)\\[4pt]
            Advancing from $L_3$ to $L_4$ yields no additional return on coordination Tiger, as strategic reasoning saturates at $L_3$.
          \end{minipage}%
        }
      \end{tabular}
      \vspace{3pt}
    \end{minipage}%
  }
  \caption{Nested initial belief construction and strategic depth saturation. (A) Combinatorial construction of $4^l$ nested initial belief instances across reasoning levels $L_1$ to $L_3$. (B) Empirical action agreement with level-$3$ policies and compute cost across $L_1$ to $L_4$ on coordination Multi-Agent Tiger, showing strategic saturation at $L_3 \to L_4$.}
  \label{fig:scaling_panel}
\end{figure}

\subsection{Convergence to Exact Bayes-Optimal Solutions on UAV Reconnaissance}

Table~\ref{tab:uav_u3} reports the performance of I-POMCP against the exact recursive I-POMDP reference solver on the Multi-Agent UAV Reconnaissance benchmark across levels $L_1, L_2, L_3$ and horizons $H \in \{2, 3, 4, 5\}$. We compare a fixed moderate simulation budget ($M_{\text{sim}} = 1{,}000$ root simulations) against horizon-scaled simulation budgets $M_{\text{sim}}^*$ that account for the exponential growth of the lookahead tree width.

\begin{table}[t]
  \caption{Evaluation of I-POMCP against Exact Reference Solvers on Multi-Agent UAV Reconnaissance across Reasoning Levels $L_1$ to $L_3$ and Horizons $H=2$ to $5$. Decision Regret ($\mathcal{L}_{\mathrm{plan}}$) and Optimal Action-Value Error (\texttt{opt\_q\_err}) are defined in Definition~\ref{def:metrics}. We compare a fixed moderate budget ($M_{\text{sim}}=1{,}000$) against horizon-adapted budgets $M_{\text{sim}}^*$.}
  \label{tab:uav_u3}
  \centering
  \small
  \resizebox{\textwidth}{!}{%
  \begin{tabular}{l c c c c c c}
    \toprule
    \textbf{Reasoning Level \& Horizon} & \textbf{Instances} & \textbf{CPU ($M_{\text{sim}}=1\text{k}$)} & \textbf{Regret ($1\text{k}$)} & \textbf{\texttt{opt\_q\_err} ($1\text{k}$)} & \textbf{Adapted Budget $M_{\text{sim}}^*$} & \textbf{Regret / Accuracy at $M_{\text{sim}}^*$} \\
    \midrule
    Level $L_1$, Horizon $H=2$ & 4 & $0.1\,\text{s}$ & $0.00$ & $0.00$ & $1{,}000$ & $0.00$ (Exact optimal) \\
    Level $L_1$, Horizon $H=3$ & 4 & $0.4\,\text{s}$ & $0.00$ & $0.00$ & $1{,}000$ & $0.00$ (Exact optimal) \\
    Level $L_1$, Horizon $H=4$ & 4 & $2.2\,\text{s}$ & $0.00$ & $1.16$ & $4{,}000$ & $0.00$ (Exact optimal) \\
    Level $L_1$, Horizon $H=5$ & 4 & $4.7\,\text{s}$ & $6.29$ & $2.40$ & $64{,}000$ & $0.00$ (Resolves $15\times$ tree width) \\
    \midrule
    Level $L_2$, Horizon $H=2$ & 16 & $1.3\,\text{s}$ & $0.00$ & $0.00$ & $1{,}000$ & $0.00$ (Exact optimal) \\
    Level $L_2$, Horizon $H=3$ & 16 & $38.0\,\text{s}$ & $0.00$ & $0.00$ & $1{,}000$ & $0.00$ (Exact optimal) \\
    Level $L_2$, Horizon $H=4$ (Exact Target) & 16 & $87.0\,\text{s}$ ($0.2\,\text{GB}$) & $0.00$ & $0.48$ & $16{,}000$ & $0.00$ ($267\text{--}269\,\text{s}$/decision) \\
    Level $L_2$, Horizon $H=4$ (Modeled Target) & 16 & N/A & N/A & N/A & $16{,}000$ & $32/32$ optimal decisions ($6{,}400\,\text{s}$) \\
    \midrule
    Level $L_3$, Horizon $H=2$ & 64 & $1.4\,\text{s}$ & $0.00$ & $0.00$ & $1{,}000$ & $0.00$ (Exact optimal) \\
    \bottomrule
  \end{tabular}%
  }
\end{table}

Three principal findings emerge from Table~\ref{tab:uav_u3}. First, even at a fixed budget of $M_{\text{sim}}=1{,}000$ simulations, I-POMCP achieves zero decision regret ($\mathcal{L}_{\mathrm{plan}} = 0.00$) across $L_1$ ($H \in \{2,3,4\}$), $L_2$ ($H \in \{2,3,4\}$), and $L_3$ ($H=2$), requiring only $87.0\,\text{s}$ total CPU time and $0.2\,\text{GB}$ peak memory across all $16$ instances at $L_2, H=4$. At $L_1, H=5$, where the $15\times$ wider branching factor causes $M_{\text{sim}}=1{,}000$ to under-sample deep observation branches ($\mathcal{L}_{\mathrm{plan}} = 6.29$), scaling the simulation budget to $M_{\text{sim}}^* = 64{,}000$ completely eliminates the approximation error and restores $\mathcal{L}_{\mathrm{plan}} = 0.00$, empirically confirming Theorem~\ref{thm:ipomcp_conv}.

Second, at $L_2, H=4$, we evaluate both an Exact-Target setting (where level-$1$ counterpart policies are queried from the exact solver) and an online Modeled-Target setting (where the level-$2$ planner solves queried level-$1$ counterpart beliefs online via nested tree search). At $M_{\text{sim}}^* = 16{,}000$, I-POMCP recovers $32/32$ ($100\%$) of the Bayes-optimal decisions across the evaluated belief instances ($6{,}400\,\text{s}$ total compute for online level-$1$ belief queries). Moreover, when solved level-$1$ counterpart policies are memoized across the search tree, increasing the root simulation budget at $L_2$ incurs negligible marginal overhead: the CPU time per decision remains virtually invariant at $267$ to $269\,\text{s}$ across root budgets of $1{,}000$, $4{,}000$, and $16{,}000$ simulations.

Third, across all $1{,}016$ evaluated reference decisions ($520$ at $L_1$, $176$ exact-target rows at $L_2$, and $320$ at $L_3$), the certified interval $[\underline{Q}(h_0, a), \overline{Q}(h_0, a)]$ contains the true optimal value $Q_H^*(b_0, a)$ in $1{,}016$ of $1{,}016$ cases ($100\%$).

\subsection{Reasoning Depth Hierarchy (\texorpdfstring{$L_0$}{L0} to \texorpdfstring{$L_4$}{L4}), Comparison with IPF+RTS, and Prior Sensitivity}

On the Multi-Agent Tiger benchmark, we evaluate a paired multi-cell experimental design across competitive and coordination payoff structures spanning reasoning levels $L_0$ through $L_4$. Across levels $L_1$ to $L_3$, I-POMCP achieves statistically significant episode return gains over the level-$0$ random walk ($\pi^0$), memoryless reactive ($\pi^{\mathrm{react}}$), and exact-belief threshold heuristic ($\pi^{\mathrm{heur}}$) policies, as well as higher mean episode returns than the particle-based $\mathrm{IPF+RTS}$ baseline \citep{doshi2009monte}, which suffers from particle depletion when informative door-creak observations falsify the majority of unweighted particles.

To investigate the marginal value of very deep recursive reasoning, we extended the coordination Tiger evaluation to level $L_4$ (expending $69\,\text{CPU-h}$ across paired trials, Figure~\ref{fig:scaling_panel}B). Interestingly, advancing from $L_3$ to $L_4$ yields no additional return gain on coordination Tiger: the level-$4$ agent selects the exact same action as the level-$3$ agent in $34$ to $35$ of $40$ paired trials ($85.0\%$ to $87.5\%$ behavioral equivalence), demonstrating that strategic reasoning depth saturates at $L_3$ in two-state coordination games. Finally, we evaluated prior sensitivity at levels $L_3$ and $L_4$ across $6$ misspecified initial belief conditions using the Holm-Bonferroni step-down multiple-testing correction \citep{holm1979simple}, confirming that canonical compressed filtering rapidly corrects misspecified base-level priors $b^{(1)}$ within $2$ to $3$ acoustic observations.

\section{Conclusion}
\label{sec:conclusion}

We have introduced Interactive POMCP (I-POMCP), an online Monte Carlo tree search algorithm for finitely nested I-POMDPs across arbitrary reasoning levels. By uniting canonical support-bounded belief filtering ($N=64$ atoms with systematic resampling) and recursive certified value bounds $[lo_r, hi_r]$ for unsampled observation branches, I-POMCP converges to exact Bayes-optimal decisions across all validated horizons on Multi-Agent Tiger ($L_0$ to $L_4$) and Multi-Agent UAV Reconnaissance ($L_1$ to $L_3$, horizons $H=2$ to $5$), achieving $1{,}016/1{,}016$ ($100\%$) empirical interval containment while outperforming particle-based Reachability Tree Sampling baselines. Future work will extend I-POMCP with learned value heuristics to further accelerate deep-horizon multi-agent search.

\begin{thebibliography}{10}

\bibitem[Doshi and Gmytrasiewicz(2009)]{doshi2009monte}
Prashant Doshi and Piotr~J. Gmytrasiewicz.
\newblock Monte {C}arlo sampling methods for approximating interactive {POMDPs}.
\newblock \emph{Journal of Artificial Intelligence Research}, 34:\penalty0 297--337, 2009.

\bibitem[Gmytrasiewicz and Doshi(2005)]{gmytrasiewicz2005framework}
Piotr~J. Gmytrasiewicz and Prashant Doshi.
\newblock A framework for sequential planning in multiagent settings.
\newblock \emph{Journal of Artificial Intelligence Research}, 24:\penalty0 49--79, 2005.

\bibitem[Holm(1979)]{holm1979simple}
Sture Holm.
\newblock A simple sequentially rejective multiple test procedure.
\newblock \emph{Scandinavian Journal of Statistics}, 6\penalty0 (2):\penalty0 65--70, 1979.

\bibitem[Kaelbling et~al.(1998)Kaelbling, Littman, and Cassandra]{kaelbling1998planning}
Leslie~Pack Kaelbling, Michael~L. Littman, and Anthony~R. Cassandra.
\newblock Planning and acting in partially observable stochastic domains.
\newblock \emph{Artificial Intelligence}, 101\penalty0 (1-2):\penalty0 99--134, 1998.

\bibitem[Kocsis and Szepesv{\'a}ri(2006)]{kocsis2006bandit}
Levente Kocsis and Csaba Szepesv{\'a}ri.
\newblock Bandit based {M}onte-{C}arlo planning.
\newblock In \emph{European Conference on Machine Learning}, pages 282--293. Springer, 2006.

\bibitem[Silver and Veness(2010)]{silver2010monte}
David Silver and Joel Veness.
\newblock Monte-{C}arlo planning in large {POMDPs}.
\newblock In \emph{Advances in Neural Information Processing Systems}, volume~23, pages 2164--2172, 2010.

\bibitem[Zeng and Doshi(2012)]{zeng2012exploiting}
Yifeng Zeng and Prashant Doshi.
\newblock Exploiting model equivalences in solving interactive dynamic influence diagrams.
\newblock \emph{Journal of Artificial Intelligence Research}, 43:\penalty0 211--255, 2012.

\end{thebibliography}

\end{document}
